\section{Formal-verification map and reproduction}\label{sec:verification}

The mathematical arguments are given in the body and Appendix~\ref{sec:topology};
reading them does not require Lean or a supplementary document. This appendix
locates the principal formal declarations and gives a complete command sequence
for checking them. Formal verification establishes the statements under the
encoded definitions; external review of those definitions and of the literature
claims remains separate.

\subsection{From paper statements to formal declarations}

The table uses the pinned [Lean] links of the paper. File paths are relative to
\texttt{formalization/}; all names are in namespace \texttt{MathResearch}.
The prefix \texttt{PC.} below abbreviates
\nolinkurl{MathResearch.PolynomialCalculus.}, and \texttt{TP.} abbreviates
\nolinkurl{MathResearch.ThirdParty.}; unprefixed names have prefix
\texttt{MathResearch.} Each file's header specifies its scope and dependencies.
The table selects the principal declarations rather than every helper lemma;
the per-statement links in the body cover the remaining steps.

\begingroup
\small
\renewcommand{\arraystretch}{1.3}
\begin{longtable}{@{}>{\raggedright\arraybackslash}p{0.30\textwidth}>{\raggedright\arraybackslash}p{0.65\textwidth}@{}}
\toprule
Paper statement and scope & Lean file and declaration\\
\midrule
\endfirsthead
\toprule
Paper statement and scope & Lean file and declaration\\
\midrule
\endhead
\bottomrule
\endfoot
Matching extension over \(\F\), Theorem~\ref{thm:moments}
& \leanref{claims/MatchingMomentExtension.lean}{\nolinkurl{claims/MatchingMomentExtension.lean}}\newline
\nolinkurl{PC.matching_moments_extend}\\
Chessboard filling over \(\F\), Corollary~\ref{cor:chessboard-filling}
& \leanref{third-party-claims/ChessboardFillingProof.lean}{\nolinkurl{third-party-claims/ChessboardFillingProof.lean}}\newline
\nolinkurl{TP.chessboard_filling}\\
Cube degree drop, Lemma~\ref{lem:cube-drop}
& \leanref{claims/BinaryCubeDegreeDrop.lean}{\nolinkurl{claims/BinaryCubeDegreeDrop.lean}}\newline
\nolinkurl{PC.cubeDifference_degree}\\
Coefficient isolation, Lemma~\ref{lem:row-cube}
& \leanref{claims/RowCubeCoefficientIsolation.lean}{\nolinkurl{claims/RowCubeCoefficientIsolation.lean}}\newline
\nolinkurl{PC.row_cube_coefficient_isolation}\\
Old-system separator, Theorem~\ref{thm:row-separation}
& \leanref{claims/CubeResidualDualSeparation.lean}{\nolinkurl{claims/CubeResidualDualSeparation.lean}}\newline
\nolinkurl{PC.cube_residual_dual_separation}\\
Bounded cofactors for ordinary restriction, Lemma~\ref{lem:restriction-ideal}
& \leanref{claims/AffineRestrictionIdeal.lean}{\nolinkurl{claims/AffineRestrictionIdeal.lean}}\newline
\nolinkurl{affine_restriction_coefficients}\\
Common restriction kernel, Lemma~\ref{lem:kernel}
& \leanref{claims/CommonAffineRestrictionKernel.lean}{\nolinkurl{claims/CommonAffineRestrictionKernel.lean}}\newline
\nolinkurl{common_affine_kernel_of_square}\\
Removal including unit-span blocks, Lemma~\ref{lem:hybrid}
& \leanref{claims/AffineFamilyRemovalWithUnits.lean}{\nolinkurl{claims/AffineFamilyRemovalWithUnits.lean}}\newline
\nolinkurl{affine_family_removal_with_units}\\
Finite exclusion with the square condition, Theorem~\ref{thm:affine}
& \leanref{claims/AffineFamilyExclusion.lean}{\nolinkurl{claims/AffineFamilyExclusion.lean}}\newline
\nolinkurl{affine_family_exclusion_of_square}\\
Actual DAG registry, Theorem~\ref{thm:dag-registry}
& \leanref{claims/AffineDAGRegistry.lean}{\nolinkurl{claims/AffineDAGRegistry.lean}}\newline
\nolinkurl{PC.affine_dag_registry}\\
Usual CNF-to-PC transfer, Theorem~\ref{thm:clause-transfer}
& \leanref{claims/BitPHPClauseTransfer.lean}{\nolinkurl{claims/BitPHPClauseTransfer.lean}}\newline
\nolinkurl{PC.usual_bitPHP_PC_transfer}\\
Main bound, both rule conventions, Theorem~\ref{thm:main}
& \leanref{claims/BitPHPExponential.lean}{\nolinkurl{claims/BitPHPExponential.lean}}\newline
\nolinkurl{bitPHP_exponential}; \nolinkurl{bitPHP_exponential_base_two}\\
Generic sufficient condition over \(\F\), Theorem~\ref{thm:generic}
& \leanref{claims/GenericCNFSubspaceCriterion.lean}{\nolinkurl{claims/GenericCNFSubspaceCriterion.lean}}\newline
\nolinkurl{PC.generic_CNF_subspace_criterion}\\
Density form, Corollary~\ref{cor:generic-density}
& \leanref{claims/GenericCNFDensityBound.lean}{\nolinkurl{claims/GenericCNFDensityBound.lean}}\newline
\nolinkurl{PC.generic_CNF_density_bound_simple}\\
\end{longtable}
\endgroup

The homological argument is proved in Lean, including the chessboard filling
used by matching extension; it is not imported as an unchecked axiom.
The final transitive axiom checks described below cover these supporting
results as well as the headline theorem.

\subsection{A pinned checkout and verification commands}

Use a machine with Git, Python 3, a C toolchain and \texttt{make}, and
\href{https://lean-lang.org/install/manual/}{elan}, Lean's toolchain manager,
on the command path. Network access is needed for the checkout, the pinned
Lean toolchain and Mathlib dependencies, and their compiled library cache.
No Python packages are required. Allow disk space for the Lean toolchain and
library build artifacts. The project pins Lean \texttt{4.34.0-rc2};
\texttt{lakefile.lean} and \texttt{lake-manifest.json} pin Mathlib and its
transitive dependencies. Do not update these pins.

The following shell commands create a fresh checkout at commit
\href{https://github.com/kbr-/math-research/tree/8904bf09a5376f48a00d1c25d079bf0c6441502a}
{\texttt{8904bf09}}, which contains the aggregate verification module and the
theorem sources used by every [Lean] link in this paper. Revision 2 changes
exposition only, so it uses the same verification target. The module's name
\texttt{BitPHPPreprintRevision1} is retained deliberately.

\begingroup\small
\begin{verbatim}
git clone https://github.com/kbr-/math-research.git bit-php-check
cd bit-php-check
git checkout --detach 8904bf09a5376f48a00d1c25d079bf0c6441502a
cd formalization
export MATHLIB_NO_CACHE_ON_UPDATE=1
export LEAN_NUM_THREADS=2
lean --version
lake exe cache get claims/BitPHPPreprintRevision1.lean
python3 verify.py --target claims/BitPHPPreprintRevision1.lean \
  --out verification.txt
lake env leanchecker --fresh claims.BitPHPPreprintRevision1
\end{verbatim}
\endgroup

On a fresh checkout the cache command retrieves the required library artifacts;
the verifier builds the project proofs and prints the types and transitive
axioms of the aggregate's listed declarations. It rejects unfinished proofs
and nonstandard axioms. The output file is written in \texttt{formalization/}; use
a new name if repeating the command. All commands must exit successfully.
The final command replays the aggregate and its imports in a fresh Lean kernel
environment. A successful build alone does not substitute for the type and
axiom checks or that replay.

For example, the expected main-theorem axiom report is
\begin{quote}\small\ttfamily
'MathResearch.bitPHP\_exponential' depends on axioms:\par
[propext, Classical.choice, Quot.sound]
\end{quote}
These are Lean's standard foundations; neither \texttt{sorryAx} nor a
custom mathematical axiom is allowed. The aggregate also checks the
superpolynomial corollary, both exponential forms, the generic criterion
and density form, and the complementary-pair short-proof control, with their
imports. It does not formalize the informal Tseitin remark, the literature
comparison, or the AI-assistance disclosure.
